\section{Conclusions}

In this work I have developed a tensor-based framework for the
analysis and transformation of centrifugal distortion parameters in
high-resolution rotational spectroscopy. The central idea is to treat
the reduced quartic and sextic tensorial objects as the primary
quantities and to interpret Watson centrifugal distortion constants as
coordinates obtained by projection onto a particular effective
Hamiltonian parametrization.

The tensorial viewpoint itself is not new, since it is already implicit
in classical rovibrational perturbation theory and in Watson's
effective Hamiltonian formalism. The novelty of the present
formulation lies instead in making this structure explicit and using it
as the basis of a practical framework that connects the direct and
inverse problems. On the direct side, the relevant tensor quantities
can be obtained from quantum-mechanical force fields and rovibrational
couplings. On the inverse side, the same tensor objects can be
reconstructed from spectroscopic Watson constants through
Moore--Penrose pseudoinverse reconstruction. In this sense the present
approach does not introduce a new phenomenological Hamiltonian, but
rather reorganizes the existing formal structure into a tensorial
representation that can be used directly for analysis and
representation transforms.

For quartic centrifugal distortion the natural reduced tensor is the
symmetric $3\times3$ matrix $\tau'$, obtained from the pair--pair
quartic tensor. This tensor admits a spectral decomposition in terms
of three eigenvalues and three orientation coordinates, thereby
providing a geometric description of quartic centrifugal distortion
prior to the choice of Watson parametrization. The passage from the
tensorial description to the five spectroscopic quartic constants can
then be viewed as a projection from the six-dimensional tensor space
onto the effective-Hamiltonian coordinate space. Conversely, the
inverse reconstruction from Watson constants becomes an affine gauge
problem, in which the Moore--Penrose pseudoinverse selects a
distinguished representative within the one-parameter family of
equivalent quartic tensors. For the standard $H_{12}H_{12}$
contribution, the Coriolis tensor and the reduced tensor $\tau'$ were
shown to encode the same spatial information in dual form, thereby
establishing a direct link between Coriolis rovibrational couplings
and quartic spectroscopic constants.

For sextic centrifugal distortion the tensor space is intrinsically
richer and does not collapse to a single symmetric second-rank tensor.
Nevertheless, the analysis developed here shows that the physical
sextic sector entering the representation-transform algorithm is not
an arbitrary numerical reduction. Operationally it appears as a
canonical $5+2$ decomposition, consisting of five physical
coordinates and a two-dimensional residual complement. Structurally,
the same sector can be interpreted as the direct sum of one isotropic
scalar and the even rank-six harmonic sector, that is, as a natural
$1+4$ decomposition. In the representation-adapted viewpoint relevant
to cyclic relabelings of the principal axes, this same
five-dimensional content acquires the form $1\oplus1\oplus1\oplus2$.
The sextic framework therefore admits a clear geometric and
group-theoretical interpretation even though it does not reduce to a
single tensor of quartic type.

These observations lead naturally to a tensor-based algorithm for
representation transforms. Starting from Watson centrifugal distortion
constants obtained either from spectroscopic fits or from
quantum-chemical calculations, one reconstructs the reduced tensorial
object, permutes the molecular axes, and finally reprojects the
transformed tensor onto Watson constants in the new representation.
For quartic centrifugal distortion this procedure reproduces the
analytical transformation formulas of Yamada and Winnewisser with
machine precision. For sextic centrifugal distortion, where analogous
representation transforms are not generally available in closed form,
the same reconstruction strategy provides a practical and numerically
stable route to internally consistent transformations.

The usefulness of the method was illustrated for dimethylsulfoxide,
for a series of HBN isomers, and for the sextic constants of
formaldehyde and thioformaldehyde. In the DMSO case the analysis shows
that large variations of individual Watson constants under
representation transforms arise primarily from the parametrization of
the effective Hamiltonian and from the numerical conditioning of the
chosen reduction--representation combination. In the HBN series the
tensor framework provides a more robust basis for comparing quartic
centrifugal distortion across related isomers, since the comparison
can be carried out at the level of tensor invariants rather than at
the level of representation-dependent fitted parameters. For the
sextic examples the decomposition into canonical physical coordinates
and residual components demonstrates that the algorithm provides not
only a representation-transform procedure but also a diagnostic tool
for assessing how closely a given sextic parameter set follows the
modeled physical subspace.

An important practical consequence of the present framework is that
representation transforms can be applied directly to spectroscopic
parameter sets without requiring explicit analytical transformation
matrices tailored to each reduction and axis system. The same approach
also provides a natural interface between spectroscopic fits,
quantum-chemical calculations, and direct rovibrational calculations.
In this way the tensor reconstruction procedure establishes a common
tensorial language for centrifugal distortion parameters obtained from
heterogeneous sources of information and makes it possible to compare
the conditioning of different representation choices before
undertaking a final refit.

A macOS application implementing the quartic and sextic
representation-transform algorithms discussed in this work is
available from the author upon request. The program allows one to
transform and analyse centrifugal distortion parameters starting from
spectroscopic Watson constants or from already computed theoretical
parameter sets, while providing tensor invariants, condition numbers,
and canonical reduced coordinates for diagnostic and interpretative
purposes.

The tensor formulation developed here should therefore be viewed not
simply as an alternative parametrization of Watson's Hamiltonian, but
as a framework that makes explicit the relation between tensorial
rovibrational content and spectroscopic effective-Hamiltonian
coordinates. This provides both a clearer conceptual picture of
centrifugal distortion and a practical computational tool for
performing representation transforms, analysing tensor invariants,
comparing spectroscopic and quantum-chemical descriptions on the same
footing, and diagnosing the numerical conditioning of alternative
spectroscopic parametrizations.
